%%% XeLaTeX-article %%%
%# -*- coding: utf-8 -*-
%!TEX encoding = UTF-8 Unicode
%!TEX TS-program = xelatex
%---------------------虽然加了%还是要保留！

\documentclass[17pt]{beamer}
\mode<presentation>
{
\usetheme[width=40pt]{Hannover}
\usecolortheme[]{dove}
\usefonttheme[]{structurebold}
\setbeameroption{hide notes}
}

\usepackage{fontspec}
\usepackage{polyglossia}
\setmainfont{Arial} %设置主字体
\newfontfamily\sanskritfont[Script=Devanagari,Mapping=romantodevanagari,Scale=1.15]{Sanskrit 2003}             %输出天城体
%\newfontfamily\sanskritfont[Mapping=tex-text]{Times New Roman}              %输出转写
\doublehyphendemerits=-10000
\newcommand{\skt}[1]{{\sanskritfont{#1}}} %输出天城体
\newcommand{\skttrans}[1]{{\skt{#1}~#1}}  %输出天城体和转写
%----------------------------------------------------设置梵文输入方法 danda । ॥

\usepackage[UTF8,fontset=windows]{ctex}
\usepackage{amsmath}
%----------------------------------------------------设置中文环境

\usepackage{graphicx}
\usepackage{flafter}
\graphicspath{{pic/}}
\usepackage{booktabs}
\usepackage{nicematrix}
\newenvironment{indentlist}
  {\begin{list}{}{\setlength{\leftmargin}{2em}\setlength{\itemsep}{0.5em}}}
  {\end{list}}
%-----------------------------------------插图表格

\usepackage{hyperref}
\usepackage[dvipsnames]{xcolor}
\usepackage{colortbl}
\definecolor{light-gray}{gray}{0.9}
%------------------------------颜色

\newcommand{\verbroot}[1]{\textcolor{red}{$\sqrt{}$#1}}
\newcommand{\sktroot}[1]{{\verbroot{\skt{#1}}}}
\newcommand{\skttransroot}[1]{{\sktroot{#1}~\textcolor{red}{#1}}}

\newcommand{\nounstem}[1]{\textcolor{red}{#1\nobreakdash-}}
\newcommand{\sktnounstem}[1]{{\textcolor{red}{\skt{#1\nobreakdash-}}}}
\newcommand{\skttransnounstem}[1]{{\sktnounstem{#1}~\nounstem{#1}}}

\newcommand{\verbstem}[1]{\textcolor{blue}{#1\nobreakdash-}}
\newcommand{\sktverbstem}[1]{{\textcolor{blue}{\skt{#1\nobreakdash-}}}}
\newcommand{\skttransverbstem}[1]{{\sktverbstem{#1}~\verbstem{#1}}}

\newcommand{\wordending}[1]{\textcolor{Orange}{\nobreakdash-#1}}
\newcommand{\sktending}[1]{{\textcolor{Orange}{\skt{-#1}}}}
\newcommand{\skttransending}[1]{{\sktending{#1}~\wordending{#1}}}

\newcommand{\fullpada}[1]{\textcolor{OliveGreen}{#1}}
\newcommand{\sktpada}[1]{{\textcolor{OliveGreen}{\skt{#1}}}}
\newcommand{\skttranspada}[1]{{\sktpada{#1}~\fullpada{#1}}}

\newcommand{\pratyaya}[1]{\mbox{\textcolor{Plum}{#1}}}
\newcommand{\sktpratyaya}[1]{{\textcolor{Plum}{\skt{#1}}}}
\newcommand{\skttranspratyaya}[1]{{\sktpratyaya{#1}~\pratyaya{#1}}}

\newcommand{\reconstruction}[1]{\textcolor{gray}{*#1}}
\newcommand{\fullsentence}[1]{\textcolor{MidnightBlue}{#1}}
\newcommand{\sktsentence}[1]{\textcolor{MidnightBlue}{\skt{#1}}}

\newcommand{\veryimportant}[1]{\textcolor{red}{#1}}
\newcommand{\important}[1]{\textcolor{blue}{#1}}
\newcommand{\notsoimportant}[1]{\textcolor{gray}{#1}}
\newcommand{\dicturl}[2]{\href{#1}{\mbox{\texttt{#2}}}}
%-------------------------------------------词根等标颜色

\title{{梵语提高}}
\subtitle{39-40. 代词五·不规则名词}
\author[张雪杉]{文学院~~张雪杉\\zhangxueshan@sdnu.edu.cn}
\date{}

\begin{document}

\begin{frame}
  \titlepage
\end{frame}

\begin{frame}
  \frametitle{本节内容}
  \small
  \tableofcontents
\end{frame}

\section{上节作业}

\begin{frame}{《薄伽梵歌》1.41--44}
  \footnotesize{
    \begin{verse}
      \skt{adharmābhibhavātkṛṣṇa praduṣyanti kulastriyaḥ ।}\\
      \mbox{\skt{strīṣu duṣṭāsu vārṣṇeya jāyate varṇasaṃkaraḥ ।। 1-41 ।।}}\\
      \skt{saṃkaro narakāyaiva kulaghnānāṃ kulasya ca ।}\\
      \mbox{\skt{patanti pitaro hyeṣāṃ luptapiṇḍodakakriyāḥ ।। 1-42 ।।}}\\
      \skt{doṣairetaiḥ kulaghnānāṃ varṇasaṃkarakārakaiḥ ।}\\
      \mbox{\skt{utsādyante jātidharmāḥ kuladharmāśca śāśvatāḥ ।। 1-43 ।।}}\\
      \skt{utsannakuladharmāṇāṃ manuṣyāṇāṃ janārdana ।}\\
      \mbox{\skt{narake 'niyataṃ vāso bhavatītyanuśuśruma ।। 1-44 ।।}}\\
    \end{verse}
  }
\end{frame}

\begin{frame}{《薄伽梵歌》1.45--47}
  \footnotesize{
    \begin{verse}
      \skt{aho bata mahatpāpaṃ kartuṃ vyavasitā vayam ।}\\
      \mbox{\skt{yadrājyasukhalobhena hantuṃ svajanamudyatāḥ ।। 1-45 ।।}}\\
      \skt{yadi māmapratīkāramaśastraṃ śastrapāṇayaḥ ।}\\
      \mbox{\skt{dhārtarāṣṭrā raṇe hanyustanme kṣemataraṃ bhavet ।। 1-46 ।।}}\\
      \skt{evamuktvārjunaḥ saṃkhye rathopastha upāviśat ।}\\
      \mbox{\skt{visṛjya saśaraṃ cāpaṃ śokasaṃvignamānasaḥ ।। 1-47 ।।}}\\
    \end{verse}
  }
\end{frame}

\begin{frame}{《薄伽梵歌》1.41--47}
  \small
~~~~~~一旦非法滋生，族中妇女堕落；一旦妇女堕落，种姓也就混乱。（41）种姓混乱导致家族和毁灭家族者堕入地狱；祖先失去供品饭和水，也跟着遭殃，纷纷坠落。​（42）这些人毁灭家族，造成种姓混乱；这些人犯下罪过，破坏宗法和种姓法。​（43）我们已经听说，折磨敌人者啊！毁弃宗法的人，注定住进地狱。​（44）由于贪图王国，贪图幸福，天哪！我们决心犯大罪，准备杀害自己人。​（45）我宁可手无寸铁，在战斗中不抵抗，让持国的儿子们，手持武器杀死我。​”​（46）阿周那说完这些话，心中充满忧伤，他放下弓和箭，坐在车座上。​（47）
\end{frame}

\section[asau/adas]{代词五：asau/adas-}
\begin{frame}{\insertsection}
  \small
  \tableofcontents[currentsection]
\end{frame}

\begin{frame}{概述}
  \small
  \begin{itemize}
    \item \nounstem{asau}/\nounstem{adas} "那个（远处的）"
    \item 与 \nounstem{ayam}/\nounstem{idam} "这个（近处的）"对比\\
    \notsoimportant{（第20章）}
  \end{itemize}
    \begin{NiceTabular}{cccc}
    \CodeBefore
      %\rectanglecolor{light-gray}{3-2}{5-2}
      %\rectanglecolor{light-gray}{7-2}{9-2}
    \Body % 7 columns
    %&   \multicolumn{3}{c}{现在时}  \\
    最近 & 近  & 远 & 最远  \\
    \fullpada{eṣaḥ/etad} & \fullpada{ayam/idam} & \fullpada{saḥ/tad} & \fullpada{asau/adas} \\
  \end{NiceTabular} 
  \begin{itemize}  
    \item 三个词干：\\
    \important{amu-} 基本词干\\
    \important{amī-} 阳性复数专用\\
    \important{amū-} 双数和阴性复数
  \end{itemize}
\end{frame}

\subsection{阳性和中性}
\begin{frame}{asau/adas- 阳性和中性变格}
  \footnotesize
  \centering
  \begin{NiceTabular}{|c|c|c|c|}[hvlines, rules/width=0.3pt, rules/color=gray]
    \Block{1-4}{\important{阳性}} & & & \\
    & \important{单数} & \important{双数} & \important{复数} \\
    主格 & \fullpada{asau} & \Block{2-1}{\fullpada{amū}} & \fullpada{amī} \\
    业格 & \fullpada{amum} &  & \fullpada{amūn} \\
    具格 & \fullpada{amunā} & \Block{3-1}{\fullpada{amūbhyām}} & \fullpada{amībhiḥ} \\
    为格 & \fullpada{amuṣmai} & & \Block{2-1}{\fullpada{amībhyaḥ}} \\
    从格 & \fullpada{amuṣmāt} & &  \\
    属格 & \fullpada{amuṣya} & \Block{2-1}{\fullpada{amuyoḥ}} & \fullpada{amīṣām} \\
    依格 & \fullpada{amuṣmin} & & \fullpada{amīṣu} \\
    \Block{1-4}{\important{中性}} & & & \\
    主/业格 & \fullpada{adaḥ} & \fullpada{amū} & \fullpada{amūni} \\
    其他 & \Block{1-3}{同阳性} & & \\
  \end{NiceTabular}
\end{frame}

\subsection{阴性}
\begin{frame}{asau/adas- 阴性变格}
  \small
  \centering
  \begin{NiceTabular}{|c|c|c|c|}[hvlines, rules/width=0.3pt, rules/color=gray]
    & \important{单数} & \important{双数} & \important{复数} \\
    主格 & \fullpada{asau} & \Block{2-1}{\fullpada{amū}} & \Block{2-1}{\fullpada{amūḥ}} \\
    业格 & \fullpada{amūm} &   &   \\
    具格 & \fullpada{amuyā} & \Block{3-1}{\fullpada{amūbhyām}} & \fullpada{amūbhiḥ} \\
    为格 & \fullpada{amuṣyai} & & \Block{2-1}{\fullpada{amūbhyaḥ}} \\
    从格 & \Block{2-1}{\fullpada{amuṣyāḥ}} & &  \\
    属格 &   & \Block{2-1}{\fullpada{amuyoḥ}} & \fullpada{amūṣām} \\
    依格 & \fullpada{amuṣyām} & & \fullpada{amūṣu} \\
  \end{NiceTabular}
\end{frame}


\section{不规则名词}
\begin{frame}{\insertsection}
  \small
  \tableofcontents[currentsection]
\end{frame}

\begin{frame}{概述}
  \small
  \begin{itemize}
    \item 四个不规则名词：\\
    \nounstem{go} (m./f.) '牛'\\
    \nounstem{dyo} (f.) '天空'\\
    \nounstem{path} (m.) '路'\\
    \nounstem{puṃs} (m.) '人'
    \item 共同特点：\important{区分强弱语干}\\
    强语干：主/呼/业格单双，主/呼格复\\
    弱语干：其他格
  \end{itemize}
\end{frame}

\subsection{go- 牛}
\begin{frame}{go- (m./f.) '牛'}
  \footnotesize
  \begin{itemize}
    \item 强语干三合：\nounstem{gau}/\nounstem{gāv}\\
    \item 弱语干二合：\nounstem{go}/\nounstem{gav}
    \item 业格单数 \fullpada{gām} 和复数 \fullpada{gāḥ} 特殊
  \end{itemize}
  \centering
  \resizebox{0.7\textwidth}{!}{%
  \begin{NiceTabular}{|c|c|c|c|}[hvlines, rules/width=0.3pt, rules/color=gray]
    & \important{单数} & \important{双数} & \important{复数} \\
    主/呼格 & \cellcolor{light-gray}\fullpada{gauḥ} & \cellcolor{light-gray}\Block{2-1}{\fullpada{gāvau}} & \cellcolor{light-gray}\fullpada{gāvaḥ} \\
    业格 & \cellcolor{light-gray}\fullpada{gām} & \cellcolor{light-gray}  & \fullpada{gāḥ} \\
    具格 & \fullpada{gavā} & \Block{3-1}{\fullpada{gobhyām}} & \fullpada{gobhiḥ} \\
    为格 & \fullpada{gave} &   & \Block{2-1}{\fullpada{gobhyaḥ}} \\
    从格 & \Block{2-1}{\fullpada{goḥ}} &   &  \\    
    属格 &  & \Block{2-1}{\fullpada{gavoḥ}} & \fullpada{gavām} \\
    依格 & \fullpada{gavi} &   & \fullpada{goṣu} \\
  \end{NiceTabular}%
  }
\end{frame}

\subsection{dyo- 天空}
\begin{frame}{dyo- (f.) '天空'}
  \footnotesize
  \begin{itemize}
    \item 有两种变化方式，第二种更常用\\
    \item 一种类似 \nounstem{go}，强语干三合，弱语干二合
    \item 一种用零级词干，元音前 \nounstem{div}，辅音前 \nounstem{dyu}
  \end{itemize}
  \centering
  \resizebox{1\textwidth}{!}{%
  \begin{NiceTabular}{|c|c|c|c|c|c|c|}[hvlines, rules/width=0.3pt, rules/color=gray]
    & \important{单数} & \important{双数} & \important{复数} & \important{单数} & \important{双数} & \important{复数} \\
    主/呼 & \cellcolor{light-gray}\fullpada{dyauḥ} & \cellcolor{light-gray}\Block{2-1}{\fullpada{dyāvau}} & \cellcolor{light-gray}\fullpada{dyāvaḥ} & \cellcolor{light-gray}\fullpada{dyauḥ} & \Block{2-1}{\fullpada{divau}} & \fullpada{divaḥ}  \\
    业格 & \cellcolor{light-gray}\fullpada{dyām} & \cellcolor{light-gray}  & \fullpada{dyūn} & \fullpada{divam} &   & \fullpada{divaḥ} \\
    具格 & \fullpada{dyavā} & \Block{3-1}{\fullpada{dyobhyām}} & \fullpada{dyobhiḥ} & \fullpada{divā} & \Block{3-1}{\fullpada{dyubhyām}} & \fullpada{dyubhiḥ}\\
    为格 & \fullpada{dyave} &   & \Block{2-1}{\fullpada{dyobhyaḥ}} & \fullpada{dive} &   & \Block{2-1}{\fullpada{dyubhyaḥ}}\\
    从格 & \Block{2-1}{\fullpada{dyoḥ}} &   &  & \Block{2-1}{\fullpada{divaḥ}} &   &  \\    
    属格 &  & \Block{2-1}{\fullpada{dyavoḥ}} & \fullpada{dyavām} &  & \Block{2-1}{\fullpada{divoḥ}} & \fullpada{divām}\\
    依格 & \fullpada{dyavi} &   & \fullpada{dyoṣu} & \fullpada{divi} &   & \fullpada{dyuṣu}\\
  \end{NiceTabular}%
  }
\end{frame}

\subsection{path- 路}
\begin{frame}{path- (m.) '路'}
  \footnotesize
  \begin{itemize}
    \item 强语干加鼻音：\nounstem{panth} （多加 \pratyaya{-ān-}）\\
    \item 弱语干无鼻音： \nounstem{path}，辅音语尾前加 \pratyaya{-i-}
  \end{itemize}
    \centering
  \resizebox{0.85\textwidth}{!}{%
  \begin{NiceTabular}{|c|c|c|c|}[hvlines, rules/width=0.3pt, rules/color=gray]
    & \important{单数} & \important{双数} & \important{复数} \\
    主/呼格 & \cellcolor{light-gray}\fullpada{panthāḥ} & \cellcolor{light-gray}\Block{2-1}{\fullpada{panthānau}} & \cellcolor{light-gray}\fullpada{panthānaḥ} \\
    业格 & \cellcolor{light-gray}\fullpada{panthānam} & \cellcolor{light-gray} & \fullpada{pathaḥ} \\
    具格 & \fullpada{pathā} & \Block{3-1}{\fullpada{pathibhyām}} & \fullpada{pathibhiḥ} \\
    为格 & \fullpada{pathe} &   & \Block{2-1}{\fullpada{pathibhyaḥ}} \\
    从格 & \Block{2-1}{\fullpada{pathaḥ}} &   &  \\
    属格 &   & \Block{2-1}{\fullpada{pathoḥ}} & \fullpada{pathām} \\
    依格 & \fullpada{pathi} &  & \fullpada{pathiṣu} \\
  \end{NiceTabular}%
  }

\end{frame}

\subsection{puṃs- 人}
\begin{frame}{puṃs- (m.) '人'}
  \footnotesize
  \begin{itemize}
    \item 强语干： \nounstem{pumāṃs}，弱语干 \nounstem{puṃs}/辅音前 \nounstem{pum}
    \item 呼格单数有两种形式 
  \end{itemize}
  \centering
  \resizebox{0.85\textwidth}{!}{%
  \begin{NiceTabular}{|c|c|c|c|}[hvlines, rules/width=0.3pt, rules/color=gray]
    & \important{单数} & \important{双数} & \important{复数} \\
    主格 & \cellcolor{light-gray}\fullpada{pumān} & \cellcolor{light-gray}\Block{3-1}{\fullpada{pumāṃsau}} & \cellcolor{light-gray}\Block{2-1}{\fullpada{pumāṃsaḥ}} \\
    呼格 & \cellcolor{light-gray}\fullpada{puman/pumaḥ} & \cellcolor{light-gray} & \cellcolor{light-gray}  \\
    业格 & \cellcolor{light-gray}\fullpada{pumāṃsam} & \cellcolor{light-gray} & \fullpada{puṃsaḥ} \\
    具格 & \fullpada{puṃsā} & \Block{3-1}{\fullpada{pumbhyām}} & \fullpada{pumbhiḥ} \\
    为格 & \fullpada{puṃse} &  & \Block{2-1}{\fullpada{pumbhyaḥ}} \\
    从格 & \Block{2-1}{\fullpada{puṃsaḥ}} &  &  \\
    属格 &   & \Block{2-1}{\fullpada{puṃsoḥ}} & \fullpada{puṃsām} \\
    依格 & \fullpada{puṃsi} &   & \fullpada{puṃsu} \\
  \end{NiceTabular}%
  }
\end{frame}




\begin{frame}{The End}
  \begin{itemize}
    \item 还想听啥？或者再见？
  \end{itemize}
\end{frame}

\end{document}
